% Proposed R2 replacement. Not a completed quantitative Results section.
% PENDING fields require a new frozen final-mask-edit execution cohort.
\subsection{Language-guided editing constructs annotation-aware tissue and cellular counterfactual states}

We first examined how language-specified interventions were represented as
explicit changes in tissue and cellular state. An instruction-only semantic
Parser separated the biological target, direction, morphology, cellular class
and clinical context while preserving negation and explicitly stated order.
The resulting request constrained primitive selection through the executable
catalog and current-mask feasibility analysis. In the existing frozen,
catalog-derived bilingual interface audit using the frozen gpt-4.1-mini
configuration, all 87 held-out instructions
(43 Chinese and 44 English) matched the scored closed-ontology semantics and
expected organ-compatible programs. This audit evaluates interface conformance,
not unrestricted clinical language or downstream mask quality. Numerical area
and cell-count budgets were specified separately from the parsed ordinal
strength labels. \textbf{[REPORT THE EXACT PARSER AND MASK-PLANNER CONFIGURATION.]}

The executable catalog distinguishes tissue-compartment geometry, local
invasive architecture and cellular composition. Its frozen scope comprises
20 open semantic primitive types and 74 organ--annotation-profile bindings.
These bindings define available actions rather than a measured execution
success rate. Boundary expansion changes the spatial extent of annotated tumour,
whereas cell-population editing changes specified nuclei within compatible
tissue support. Similarly, changing an annotated immune-rich compartment is
distinct from adding inflammatory cells while preserving tissue labels.
Cellular construction combines source-derived population constraints,
ProbNet spatial scores and complete reference nucleus shapes.
\textbf{[PENDING: VERSION-MATCHED MULTISCALE EXECUTION COUNTS AND MEASUREMENTS.]}

Annotation profiles constrain how a biological action can be represented.
For example, complete gland annotations and pattern-specific epithelial labels
require different interpretations of gland boundaries and lumen-associated
space. The current masks and provenance-bound auxiliary structures determine
which catalog actions can be instantiated in each patch. This prevents a shared
semantic instruction from silently imposing the same geometric rule on
incompatible annotation systems.
\textbf{[PENDING: MATCHED GLaS/PANDA EXAMPLES; ADD ABLATION ONLY IF CLAIMING A
QUANTITATIVE BENEFIT OVER A SHARED RULE.]}

Candidate edits are committed only after the prescribed mask constraints pass,
including legal label transitions, requested direction and magnitude,
complete-instance placement and preservation of protected state. The joint
change region measures altered tissue or nucleus labels, whereas generation
support specifies the region available for downstream image reconstruction.
We distinguish the validity of accepted outputs from the fraction of attempted
requests that produce an accepted state, retaining clarification, abstention,
timeout and runtime failure in their appropriate denominators.
\textbf{[PENDING: FROZEN ATTEMPT-LEVEL STATUS COUNTS, TARGET ERROR AND PRESERVATION
MEASUREMENTS; DO NOT SUBSTITUTE THE LEGACY TISSUE-ONLY BENCHMARK.]}

Explicitly ordered programs update the masks after each validated step and
rebuild the current state before the next operation. In the illustrated repeated
expansion example, tumour occupancy increased from 49.3\% in the original mask
to 54.7\% after one step and 60.2\% after two steps. Each step used a 6\% joint-area
budget; these budgets are distinct from the measured final tumour-area change.
The cumulative generation interface retains the original reference, final target
masks and the union of both validated support regions (23.7\% of the patch).
Global reconstruction was explicitly selected for this example. This demonstrates
one audited repeated intervention, not arbitrary composition, automatic
large-budget decomposition or prediction of longitudinal disease progression.
\textbf{[PENDING: ADD THE PRESPECIFIED MULTI-OPERATION PROGRAM COHORT BEFORE MAKING
A GENERAL COMPOSITIONALITY CLAIM.]}
